\begin{abstract}
Modular LLM agents that route between neural reasoning and deterministic computation promise improved reliability in rule-governed environments, yet understanding \emph{what} drives effective routing remains unclear. Per-step oracle labeling degenerates in sparse-reward settings---in ScienceWorld, 98\% of steps yield tied outcomes, rendering step-level analysis uninformative. We propose episode-level counterfactual analysis: running complete episodes under pure strategies and comparing final outcomes to derive per-episode routing labels. Applying this methodology to 245 episodes across 11 task types with a Qwen2.5-7B-Instruct backbone, we find that death avoidance accounts for 90.41\% of the routing signal (Death Avoidance Fraction). Per-task death rate differentials correlate near-perfectly with routing benefit ($r$ = 0.9698, $p$ = 8.12$\times$10$^{-7}$). The oracle router agrees with always-deterministic on 96.2\% of decisions, and a trained router achieves only 58.1\% balanced accuracy---near random. These results indicate that deterministic computation's value is concentrated in catastrophic failure prevention rather than positive task guidance.
\end{abstract}
